%%===================================================================
%% Lecture 10 -- Product Groups, Accidental Symmetries, and Local Gauge Invariance
%% 77609 -- Introduction to Elementary Particles
%% Date: 10/05/2026
%% Lecturer: Prof. Yonit Hochberg
%%===================================================================

\documentclass[../main.tex]{subfiles}
\usepackage{preamble}

\begin{document}

\section{Lecture 10 -- Product Groups, Accidental Symmetries, and Local Gauge Invariance}\label{sec:lec10}
\begin{center}
\begin{tabular}{ll}
  \textbf{Date:}\quad 10/05/2026 & \\
\end{tabular}
\end{center}
\hrule\vspace{1.5em}

%%------------------------------------------------------------------
\subsection*{Overview}

This lecture extends the model-building framework in two directions.  First, we allow multiple independent $U(1)$ symmetries acting on multiple fields and discover that the renormalizability constraint can accidentally produce \emph{more} symmetry than was imposed.  Second, we promote the transformation parameter $\theta$ from a constant to a function of spacetime, arriving at the concept of a \textbf{local} (gauge) symmetry, the covariant derivative, and the necessity and properties of massless gauge bosons.

\begin{enumerate}
  \item \textbf{Product groups}: multiple independent $U(1)$ symmetries and how to assign charges under each.
  \item \textbf{Example}: two complex scalars each charged under a different $U(1)$.
  \item \textbf{Accidental symmetries}: truncating at dimension four can accidentally give a larger symmetry.
  \item \textbf{Global versus local}: what changes when $\theta\to\theta(x)$.
  \item \textbf{Covariant derivative}: how to restore invariance of kinetic terms under local $U(1)$.
  \item \textbf{Gauge field}: its transformation law and kinetic term.
  \item \textbf{Masslessness of gauge bosons}: why local gauge invariance forbids a mass term.
  \item \textbf{Multiple gauge factors}: each independent $U(1)$ factor requires its own gauge field.
\end{enumerate}

%%------------------------------------------------------------------
\subsection*{Product Groups}

\subsubsection*{Physical Motivation}

In Lecture~09, all fields rotated under a single parameter $\theta$.  Nature, however, imposes several independent phase rotations: the Standard Model gauge group $SU(3)_c\times SU(2)_L\times U(1)_Y$ already contains multiple independent factors.  We therefore need the language to handle several commuting $U(1)$ symmetries simultaneously.

A \textbf{product group} $U(1)_A\times U(1)_B\times\cdots$ is a collection of independent symmetries, each with its own transformation parameter.  Every field carries a separate charge under each factor.

\subsubsection*{Two fields, two independent $U(1)$s}

Take two complex scalar fields $\phi_1$ and $\phi_2$, and require invariance under two independent phase rotations parametrised by $\theta_A$ and $\theta_B$.  The transformations are:
\begin{equation}
  \boxed{
  \phi_i \;\longrightarrow\;
  e^{i Q_{iA}\,\theta_A + i Q_{iB}\,\theta_B}\,\phi_i,
  \qquad i=1,2,}
  \label{eq:lec10:product-group-transformation}
\end{equation}
where $Q_{iA}$ is the charge of $\phi_i$ under $U(1)_A$ and $Q_{iB}$ is the charge under $U(1)_B$.

\textit{Physical significance:} In contrast to a single $U(1)$, here there are two independent angles.  One \emph{cannot} in general write $Q_{iA}\theta_A + Q_{iB}\theta_B$ as a single angle times a single charge for both fields simultaneously, unless the charge ratios happen to coincide.

\subsubsection*{A concrete charge assignment}

Choose:
\begin{equation}
  \boxed{
  \begin{array}{c|cc}
       & Q_{A} & Q_{B} \\ \hline
  \phi_1 & 1 & 0 \\
  \phi_2 & 0 & 1
  \end{array}}
  \label{eq:lec10:charge-assignment}
\end{equation}
meaning $\phi_1$ rotates only under $U(1)_A$ and $\phi_2$ rotates only under $U(1)_B$:
\begin{equation}
  \phi_1\to e^{i\theta_A}\phi_1, \qquad
  \phi_2\to e^{i\theta_B}\phi_2.
  \label{eq:lec10:example-transformations}
\end{equation}

\subsubsection*{Most general renormalizable Lagrangian}

A term is allowed if and only if its total charge under \emph{each} $U(1)$ factor is zero.  With the charge assignment above:
\begin{itemize}
  \item Under $U(1)_A$: only $\phi_1^\dagger\phi_1$ combinations are invariant; $\phi_2$ is neutral.
  \item Under $U(1)_B$: only $\phi_2^\dagger\phi_2$ combinations are invariant; $\phi_1$ is neutral.
\end{itemize}

The most general renormalizable Lagrangian is:
\begin{equation}
  \boxed{
  \begin{aligned}
  \lagr_{\rm cocoa}
  ={}&
  \partial_\mu\phi_1^\dagger\partial^\mu\phi_1
  +\partial_\mu\phi_2^\dagger\partial^\mu\phi_2
  -m_1^2\phi_1^\dagger\phi_1
  -m_2^2\phi_2^\dagger\phi_2 \\
  &-\lambda_{11}(\phi_1^\dagger\phi_1)^2
  -\lambda_{22}(\phi_2^\dagger\phi_2)^2
  -\lambda_{12}(\phi_1^\dagger\phi_1)(\phi_2^\dagger\phi_2).
  \end{aligned}}
  \label{eq:lec10:cocoa-lagrangian}
\end{equation}

Crucially, a mixed term such as $\phi_1^3\phi_2^\dagger$ from Lecture~09 is \emph{forbidden} here: its charge under $U(1)_A$ is $3\cdot1-0=3\neq0$.  The charge assignment dictates which interactions may appear.

\textit{Physical significance:} The allowed Lagrangian is more restrictive here than in the $q_1=1,q_2=3$ example of Lecture~09, because the two fields do not talk to each other through the symmetry.  Their only coupling is the portal term $\lambda_{12}(\phi_1^\dagger\phi_1)(\phi_2^\dagger\phi_2)$, which is charge-neutral under both $U(1)$s.

%%------------------------------------------------------------------
\subsection*{Accidental Symmetries}

\subsubsection*{Physical Motivation}

The Lagrangian $\lagr_{\rm cocoa}$ was derived by imposing $U(1)_A\times U(1)_B$.  A natural question is: does this Lagrangian have more symmetry than we intended?  The answer is yes, and understanding why matters for the Standard Model.

\subsubsection*{Obtaining the same Lagrangian from a single $U(1)$}

Consider a single $U(1)$ with charges:
\begin{equation}
  \phi_1\to e^{i\cdot1\cdot\theta}\phi_1, \qquad
  \phi_2\to e^{i\cdot4\cdot\theta}\phi_2.
  \label{eq:lec10:single-u1-example}
\end{equation}
A potential mixed term would need total charge zero.  The only renormalizable candidate of this form would involve four fields; for example $\phi_1^{-4}\phi_2$ has dimension $>4$.  Hence, at the renormalizable level, the most general Lagrangian consistent with this single $U(1)$ also contains no mixed $\phi_1$-$\phi_2$ interaction terms, and one again obtains $\lagr_{\rm cocoa}$.

But $\lagr_{\rm cocoa}$ clearly has two independent phase symmetries: rotating $\phi_1$ alone or $\phi_2$ alone both leave it invariant.  So:
\begin{equation}
  \boxed{
  \text{imposed: }U(1)
  \quad\xrightarrow{\text{renormalizability}}\quad
  \lagr_{\rm cocoa}
  \quad\xleftarrow{\text{symmetry of }\lagr}\quad
  U(1)_A\times U(1)_B.}
  \label{eq:lec10:accidental-flow}
\end{equation}

\dfn{Accidental Symmetry}{An \textbf{accidental symmetry} is a symmetry of the renormalizable Lagrangian that was \emph{not} explicitly imposed.  It arises because all symmetry-breaking operators consistent with the imposed symmetry happen to have mass dimension $>4$ and are therefore absent from the renormalizable theory.}

\textit{Physical significance:} Accidental symmetries are broken by non-renormalizable operators suppressed by a high scale $\Lambda$.  Famous examples in the Standard Model are \textbf{baryon number} $B$ and \textbf{lepton number} $L$: no renormalizable SM operator violates them, so they appear exact at collider energies.  However, dimension-six operators (suppressed by $1/\Lambda^2$) do violate $B$ and $L$, predicting (unobserved so far) proton decay.

\subsubsection*{Choice of basis for the symmetry generators}

The two independent symmetries $\theta_A$ and $\theta_B$ can equally well be described by any two independent linear combinations.  A common choice is:
\begin{equation}
  \theta_+ = \theta_A + \theta_B, \qquad \theta_- = \theta_A - \theta_B,
  \label{eq:lec10:linear-combinations}
\end{equation}
giving $U(1)_+\times U(1)_-$.  Any two independent linear combinations are physically equivalent.

%%------------------------------------------------------------------
\subsection*{From Global to Local Symmetry}

\subsubsection*{Physical Motivation}

All symmetries discussed so far have been \textbf{global}: the parameter $\theta$ is the same everywhere in space and time.  This is physically strange.  Why should a rotation performed ``here'' be rigidly correlated with one performed at a distant galaxy?

A \textbf{local} (or \textbf{gauge}) symmetry removes this rigidity by allowing $\theta$ to vary independently at each spacetime point.  The price is that we must introduce a new dynamical field—the gauge field—to maintain invariance of the kinetic terms.  Gauge symmetry is the organizing principle of the Standard Model.

\subsubsection*{Definition}

\dfn{Local $U(1)$ symmetry}{A transformation
\begin{equation}
  \boxed{
  \phi(x)\;\longrightarrow\; e^{iq\,\theta(x)}\,\phi(x),}
  \label{eq:lec10:local-u1}
\end{equation}
where $\theta(x)\equiv\theta(x^\mu)$ is an arbitrary smooth function of spacetime, is called a \textbf{local} or \textbf{gauge} $U(1)$ transformation.}

\textit{Physical significance:} Global means $\partial_\mu\theta=0$.  Local means $\theta$ can take any value at any spacetime point; it is not a rigid phase.

\subsubsection*{Non-derivative terms are automatically gauge invariant}

Any term in the Lagrangian that does not involve a spacetime derivative—mass terms, potential terms—transforms as:
\begin{equation}
  \phi^\dagger\phi \;\longrightarrow\; (e^{-iq\theta(x)}\phi^\dagger)(e^{iq\theta(x)}\phi)
  = \phi^\dagger\phi.
  \label{eq:lec10:no-derivative-invariance}
\end{equation}
The phases cancel regardless of whether $\theta$ depends on $x$.  Therefore:
\begin{equation}
  \boxed{
  \text{if a non-derivative term is globally invariant, it is also locally invariant.}}
  \label{eq:lec10:nonderiv-always-invariant}
\end{equation}

\textit{Physical significance:} Only kinetic terms need special treatment when going from global to local symmetry.

%%------------------------------------------------------------------
\subsection*{The Problem with Kinetic Terms}

\subsubsection*{Physical Motivation}

The derivative $\partial_\mu$ does not commute with a position-dependent phase.  When $\theta$ varies in space, differentiating $e^{iq\theta(x)}\phi(x)$ produces an extra term involving $\partial_\mu\theta$.  This extra term breaks gauge invariance of the kinetic Lagrangian and must be corrected.

\subsubsection*{How $\partial_\mu\phi$ transforms}

Under the local transformation \eqref{eq:lec10:local-u1}:
\begin{align}
  \partial_\mu\phi \;\longrightarrow\;
  &\partial_\mu\!\left(e^{iq\theta(x)}\phi(x)\right) \nonumber \\
  ={}&
  e^{iq\theta(x)}\,\partial_\mu\phi(x)
  \;+\;
  iq\,(\partial_\mu\theta)\,e^{iq\theta(x)}\,\phi(x).
  \label{eq:lec10:derivative-transforms}
\end{align}

The first piece is the ``good'' term (what we would want from a global transformation); the second is an extra piece proportional to $\partial_\mu\theta$.  Therefore:
\begin{equation}
  \boxed{
  \partial_\mu\phi\;\longrightarrow\;
  e^{iq\theta(x)}\!\left(\partial_\mu\phi + iq(\partial_\mu\theta)\phi\right).}
  \label{eq:lec10:derivative-bad}
\end{equation}

The kinetic term $(\partial_\mu\phi^\dagger)(\partial^\mu\phi)$ is \emph{not} invariant under the local transformation because of the cross terms involving $\partial_\mu\theta$.

\textit{Physical significance:} This breakdown of kinetic-term invariance is universal: it happens for scalar and fermion kinetic terms alike.  We must repair it.

%%------------------------------------------------------------------
\subsection*{The Covariant Derivative}

\subsubsection*{Physical Motivation}

The fix is to replace the ordinary derivative $\partial_\mu$ with a \textbf{covariant derivative} $D_\mu$ that, when acting on a charged field, transforms the same way under local $U(1)$ as $\partial_\mu$ would under a global $U(1)$:
\begin{equation}
  \text{want: }\quad
  D_\mu\phi \;\longrightarrow\; e^{iq\theta(x)}\,D_\mu\phi.
  \label{eq:lec10:cov-deriv-requirement}
\end{equation}
If $D_\mu\phi$ transforms this cleanly, then $(D_\mu\phi)^\dagger(D^\mu\phi)$ is gauge invariant by the same argument as $\phi^\dagger\phi$.

\subsubsection*{Construction of the covariant derivative}

\dfn{Covariant derivative}{Introduce a new vector field $A_\mu(x)$ and define:
\begin{equation}
  \boxed{
  D_\mu \;\equiv\; \partial_\mu + igqA_\mu,}
  \label{eq:lec10:covariant-derivative}
\end{equation}
where $g$ is a dimensionless coupling constant (the \textbf{gauge coupling}), $q$ is the charge of the field being acted upon, and $A_\mu$ is the \textbf{gauge field}.}

\subsubsection*{Transformation law of the gauge field}

For $D_\mu\phi$ to satisfy requirement \eqref{eq:lec10:cov-deriv-requirement}, the gauge field must transform as:
\begin{equation}
  \boxed{
  A_\mu(x) \;\longrightarrow\; A_\mu(x) - \frac{1}{g}\,\partial_\mu\theta(x).}
  \label{eq:lec10:gauge-field-transformation}
\end{equation}

\textit{Verification:} Under simultaneous transformations \eqref{eq:lec10:local-u1} and \eqref{eq:lec10:gauge-field-transformation}:
\begin{align}
  D_\mu\phi &= (\partial_\mu + igqA_\mu)\phi \nonumber \\
  &\longrightarrow
  \left(\partial_\mu + igq\!\left(A_\mu - \tfrac{1}{g}\partial_\mu\theta\right)\right)
  e^{iq\theta}\phi \nonumber \\
  &= e^{iq\theta}\partial_\mu\phi
    + iq(\partial_\mu\theta)e^{iq\theta}\phi
    + igqA_\mu e^{iq\theta}\phi
    - iq(\partial_\mu\theta)e^{iq\theta}\phi \nonumber \\
  &= e^{iq\theta}(\partial_\mu + igqA_\mu)\phi
  = e^{iq\theta}\,D_\mu\phi. \quad\checkmark
  \label{eq:lec10:cov-deriv-check}
\end{align}

\textit{Physical significance:} The extra term $-\frac{1}{g}\partial_\mu\theta$ in the gauge field transformation precisely cancels the unwanted $\partial_\mu\theta$ term from the derivative.  This is the mechanism by which the gauge field ``absorbs'' the local phase ambiguity.

\textit{Physical significance:} The transformation law of $A_\mu$ is \textbf{additive} (a shift), not multiplicative.  It resembles a phase shift rather than a rotation.  This is a distinctive feature of abelian gauge fields.

\subsubsection*{Gauge-invariant Lagrangian for a charged scalar}

Replacing $\partial_\mu\to D_\mu$ everywhere in the scalar kinetic term:
\begin{equation}
  \boxed{
  \lagr_{\rm scalar}
  =
  (D_\mu\phi)^\dagger(D^\mu\phi) - m^2\phi^\dagger\phi - \lambda(\phi^\dagger\phi)^2,}
  \label{eq:lec10:scalar-gauge-lagrangian}
\end{equation}
which is invariant under the local $U(1)$ by construction.  The same replacement works for fermion kinetic terms: $\bar\psi i\slashed{\partial}\psi \to \bar\psi i\slashed{D}\psi$.

%%------------------------------------------------------------------
\subsection*{Kinetic Term for the Gauge Field}

\subsubsection*{Physical Motivation}

We have introduced a new dynamical field $A_\mu$.  Since it is dynamical, it must have a kinetic term in the Lagrangian.  This term must be gauge invariant by itself.  The ordinary combination $(\partial_\mu A_\nu)(\partial^\mu A^\nu)$ is not gauge invariant because $A_\mu$ shifts by a gradient under \eqref{eq:lec10:gauge-field-transformation}.  We need a combination that cancels these gradients.

\subsubsection*{Field strength tensor}

Define:
\begin{equation}
  \boxed{
  F_{\mu\nu} \;\equiv\; \partial_\mu A_\nu - \partial_\nu A_\mu.}
  \label{eq:lec10:field-strength}
\end{equation}

Under the gauge transformation \eqref{eq:lec10:gauge-field-transformation}:
\begin{equation}
  F_{\mu\nu}
  \;\longrightarrow\;
  \partial_\mu\!\left(A_\nu - \tfrac{1}{g}\partial_\nu\theta\right)
  - \partial_\nu\!\left(A_\mu - \tfrac{1}{g}\partial_\mu\theta\right)
  = F_{\mu\nu} - \tfrac{1}{g}(\partial_\mu\partial_\nu\theta - \partial_\nu\partial_\mu\theta)
  = F_{\mu\nu},
  \label{eq:lec10:fmunu-invariant}
\end{equation}
using the commutativity of partial derivatives.  Hence $F_{\mu\nu}$ is gauge invariant.

\textit{Physical significance:} $F_{\mu\nu}$ is the electromagnetic field strength.  Its components are the electric and magnetic fields.  Gauge invariance ensures that physical observables—which depend on $F_{\mu\nu}$, not on $A_\mu$ directly—are unambiguous.

\subsubsection*{Gauge-invariant kinetic Lagrangian for $A_\mu$}

The unique renormalizable (dimension-four) Lorentz scalar built from $F_{\mu\nu}$ is:
\begin{equation}
  \boxed{
  \lagr_{\rm gauge}
  =
  -\frac{1}{4}\,F_{\mu\nu}F^{\mu\nu}.}
  \label{eq:lec10:gauge-kinetic}
\end{equation}

The factor $-1/4$ is a conventional normalization (ensures a positive-definite Hamiltonian with canonically normalized fields).  The equation of motion for $A_\mu$ derived from $\lagr_{\rm gauge}$ alone is:
\begin{equation}
  \partial^\mu F_{\mu\nu} = 0,
  \label{eq:lec10:maxwell-vacuum}
\end{equation}
which is precisely Maxwell's equations in vacuum.

\textit{Physical significance:} The entire structure of electromagnetism—Maxwell's equations, the photon, gauge invariance—follows from demanding a local $U(1)$ symmetry.  The gauge field $A_\mu$ \emph{is} the photon field.

%%------------------------------------------------------------------
\subsection*{Masslessness of Gauge Bosons}

\subsubsection*{Physical Motivation}

$A_\mu$ has a kinetic term, so its quantum excitations are particles.  What is their mass?  A mass term for a vector field would normally take the form $m^2 A_\mu A^\mu$.  But gauge invariance forbids this, as we now show.

\subsubsection*{The mass term violates gauge invariance}

Under $A_\mu\to A_\mu - \frac{1}{g}\partial_\mu\theta$:
\begin{equation}
  m^2 A_\mu A^\mu
  \;\longrightarrow\;
  m^2\!\left(A_\mu - \tfrac{1}{g}\partial_\mu\theta\right)\!\left(A^\mu - \tfrac{1}{g}\partial^\mu\theta\right)
  =
  m^2 A_\mu A^\mu
  - \frac{2m^2}{g}\,A^\mu\partial_\mu\theta
  + \frac{m^2}{g^2}(\partial_\mu\theta)^2
  \;\neq\; m^2 A_\mu A^\mu.
  \label{eq:lec10:mass-term-not-invariant}
\end{equation}

Since this is not invariant, the mass term is \emph{not allowed} in the gauge-invariant Lagrangian.  Therefore:
\begin{equation}
  \boxed{
  \text{local gauge invariance}
  \;\Longrightarrow\;
  \text{gauge bosons are massless.}}
  \label{eq:lec10:masslessness}
\end{equation}

\textit{Physical significance:} The photon's masslessness is not a measured coincidence.  It is a logical consequence of requiring local $U(1)$ gauge invariance.  The photon, as the $U(1)_{\rm em}$ gauge boson, must be exactly massless.  A nonzero mass would signal the breaking of gauge invariance.

\subsubsection*{Degrees of freedom}

A massive vector boson has three polarization states; a massless vector boson has only two (transverse polarizations).  The masslessness of gauge bosons is consistent with the masslessness of the photon, which is observed to have exactly two polarizations.

%%------------------------------------------------------------------
\subsection*{Full Gauge-Invariant Lagrangian}

Combining the kinetic term for $A_\mu$ with the gauge-invariant matter Lagrangian gives the complete theory:
\begin{equation}
  \boxed{
  \lagr_{\QED}
  =
  -\frac{1}{4}F_{\mu\nu}F^{\mu\nu}
  +
  (D_\mu\phi)^\dagger(D^\mu\phi)
  -m^2\phi^\dagger\phi
  -\lambda(\phi^\dagger\phi)^2,}
  \label{eq:lec10:full-qed-scalar}
\end{equation}
with $D_\mu=\partial_\mu+igqA_\mu$, $F_{\mu\nu}=\partial_\mu A_\nu-\partial_\nu A_\mu$.

Expanding $(D_\mu\phi)^\dagger(D^\mu\phi)$ produces:
\begin{equation}
  (D_\mu\phi)^\dagger(D^\mu\phi)
  =
  (\partial_\mu\phi^\dagger)(\partial^\mu\phi)
  + igqA^\mu(\phi^\dagger\partial_\mu\phi - \phi\partial_\mu\phi^\dagger)
  + g^2q^2A_\mu A^\mu\,\phi^\dagger\phi,
  \label{eq:lec10:expanded-kinetic}
\end{equation}
showing interaction vertices with one or two gauge bosons.

\textit{Physical significance:} The single-photon vertex (proportional to $gq$) shows that charge determines the coupling strength of the field to the photon.  This is the second role of charge: it governs the interaction rate, not just the phase rotation.

\textit{Physical significance:} The $U(1)$ gauge field does \emph{not} appear in its own kinetic term with a self-coupling: $F_{\mu\nu}F^{\mu\nu}$ contains no $A^4$ terms.  The photon does not couple to itself.  This is a special property of abelian gauge groups; non-abelian gauge fields (gluons, $W$, $Z$) do self-interact.

%%------------------------------------------------------------------
\subsection*{Product Gauge Groups}

\subsubsection*{Physical Motivation}

The Standard Model gauge group is $SU(3)_c\times SU(2)_L\times U(1)_Y$.  The $U(1)$ factor alone demonstrates the logic: if one has a product gauge group $U(1)_A\times U(1)_B\times\cdots$, each factor is an independent local symmetry requiring its own treatment.

\subsubsection*{Each factor requires its own gauge field and coupling}

If the local symmetry decomposes into $n$ independent commuting $U(1)$ factors, then for each factor $\alpha=A,B,\ldots$:
\begin{itemize}
  \item There is an independent parameter $\theta_\alpha(x)$.
  \item There is an independent gauge field $A_\mu^\alpha$ with transformation law
    \[A_\mu^\alpha\to A_\mu^\alpha - \frac{1}{g_\alpha}\partial_\mu\theta_\alpha.\]
  \item There is an independent dimensionless coupling $g_\alpha$.
  \item The covariant derivative acting on a field with charges $(Q_A, Q_B,\ldots)$ is:
    \[D_\mu = \partial_\mu + ig_A Q_A A_\mu^A + ig_B Q_B A_\mu^B + \cdots\]
\end{itemize}

\begin{equation}
  \boxed{
  U(1)_A\times U(1)_B\times\cdots
  \;\Longrightarrow\;
  n\text{ massless gauge bosons},\quad
  n\text{ independent couplings }g_\alpha.}
  \label{eq:lec10:product-gauge-count}
\end{equation}

\textit{Physical significance:} The coupling constants of different gauge factors are independent: gauge symmetry does not predict the ratio $g_A/g_B$.  Each must be measured.  This is why the Standard Model has three independent gauge couplings $g_3,g_2,g_1$ that are not predicted by gauge symmetry alone—they must be measured from experiment.

%%------------------------------------------------------------------
\subsection*{Summary Table: Global vs.\ Local $U(1)$}

\begin{center}
\begin{tabular}{lll}
\toprule
 & \textbf{Global $U(1)$} & \textbf{Local $U(1)$} \\
\midrule
Parameter & $\theta=\text{const}$ & $\theta=\theta(x)$ \\
Non-derivative terms & Invariant & Invariant \\
Kinetic terms & Invariant & Not invariant (need $D_\mu$) \\
New field required & No & Yes: $A_\mu$ \\
Gauge boson mass & --- & Must be zero \\
Consequence & Conserved charge & Gauge interaction \\
\bottomrule
\end{tabular}
\end{center}

%%------------------------------------------------------------------
\subsection*{To Remember}

\begin{itemize}
  \item A product group $U(1)^n$ means $n$ independent phase rotations; each field carries $n$ charges.
  \item Renormalizability (dimension $\leq4$) can produce more symmetry than imposed: this is an \textbf{accidental symmetry}.
  \item Accidental symmetries of the SM include baryon number $B$ and lepton number $L$; they are broken by dimension-six operators suppressed by a high scale.
  \item Going from global to local: only derivative terms (kinetic terms) are affected.
  \item The covariant derivative $D_\mu = \partial_\mu + igqA_\mu$ restores gauge invariance; $A_\mu$ must transform as $A_\mu\to A_\mu - \frac{1}{g}\partial_\mu\theta$.
  \item The field strength $F_{\mu\nu}=\partial_\mu A_\nu-\partial_\nu A_\mu$ is gauge invariant; its equation of motion gives Maxwell's equations.
  \item The kinetic Lagrangian for the gauge field is $-\frac{1}{4}F_{\mu\nu}F^{\mu\nu}$.
  \item A mass term $m^2 A_\mu A^\mu$ is not gauge invariant, so gauge bosons must be \textbf{massless}.
  \item A product gauge group $U(1)^n$ requires $n$ independent gauge fields and $n$ independent couplings.
  \item The coupling $gq$ in $D_\mu\phi = (\partial_\mu + igqA_\mu)\phi$ determines how strongly the field interacts with the gauge boson.
\end{itemize}

%%------------------------------------------------------------------
\subsection*{Recommended Reading}

\begin{itemize}
  \item \textbf{Tong, \textit{Lectures on Particle Physics}, Chapter~2}
    (\href{https://www.damtp.cam.ac.uk/user/tong/pp.html}{\texttt{damtp.cam.ac.uk/user/tong/pp.html}}):
    global and local symmetries, covariant derivatives, gauge fields.
  \item \textbf{Peskin \& Schroeder, Section~4.1}: gauge invariance, covariant derivative, $F_{\mu\nu}$.
  \item \textbf{Griffiths, \textit{Introduction to Elementary Particles}, Chapter~9}: gauge invariance and the photon as the gauge boson of $U(1)_{\rm em}$.
  \item \textbf{Halzen \& Martin, Chapter~14}: gauge theories, photon coupling, and QED.
\end{itemize}

\end{document}
